\documentclass[12 pt]{exam}
\usepackage[margin=0.75in]{geometry}
\usepackage{amsmath,amssymb, amsthm, color}
\usepackage{enumerate,graphicx,multicol}
\usepackage{wrapfig}
\usepackage{pgf,tikz, subcaption}
\usetikzlibrary{arrows}

\noprintanswers
\newtheorem{definition}{Definition}



\pagestyle{head}                       % This causes the exam to only have headers, no footers.
%\runningheadrule                     % This causes the headings to have an underline.

%\firstpageheader{Name:\enspace\hbox to 4.5in{\hrulefill} Section:\enspace\hbox to 1.2in{\hrulefill}}{}{}
%\extraheadheight{.25in}

\begin{document}
\hrule
\begin{center}
  \huge{Math 2710 Problem Set 1: \textsection 1.1-1.4}    %Title
\end{center}
\hrule
\vskip .2in


\begin{questions}

\question For each of the following sets, (i) write the set in set builder notation, i.e. $\left\{x : p(x) \right\}$. Then (ii) write the set as a list of its elements, if possible:
	\begin{parts}
		\part The set of natural numbers whose square is less than 25
		\part $(-1,8]$
		\part The set of rational numbers less than 4
	\end{parts}
 
\begin{solution}
Enter your solution here.
\end{solution} 

\question Determine whether each of the following is true or false. Explain your reasoning for each answer:
	\begin{parts}
		\part The empty set is a proper subset of every set.
		\part The empty set is an element of every set.
		
	%\end{parts}
	
%\question Determine whether each of the following is true or false. Explain your reasoning for each answer:
%	\begin{parts}
		\part	$\{1\} \in \{1,2,3,\{4\}\}$
		\part $\{1\} \subseteq \{1,2,3,\{4\}\}$
		
%	\end{parts}
	
%	\question Determine whether each of the following is true or false. Explain your reasoning for each answer:
%	\begin{parts}
		\part $\{4\} \subseteq \{1,2,3,\{4\}\}$
		\part $\{3\} \subset \mathcal{P}(\mathbb{Q})$
		\part $\{\{\emptyset\}\} \subseteq \mathcal{P}(\{\emptyset, \{\emptyset\}\})$
		
	\end{parts}
	
\question Give an example, if possible, of sets $A,B,$ and $C$ satisfying the following. If it is not possible, explain why not.
	\begin{parts}
		\part $A\subseteq B$, $B \subseteq C$, and $C \subseteq A$
		\part $A\subset B$, $B \subset C$, and $C \subset A$
		\part $A \subseteq B$, $B \nsubseteq C$, and $C \nsubseteq A$
	\end{parts}

\question Write the power set of each of the following sets.
	\begin{parts}
		\part $\{\emptyset, \triangle, \square\}$
		\part $\{a, \{b,\{c\} \} \}$
	\end{parts}
	

\question For $A=\{ 1, 2\}$, determine $A \times \mathcal{P}(A)$.

\question Assume $A$ and $B$ are sets with $|A|=4$ and $|B|=3.$ Determine $|A \times B|$. Explain your reasoning. 

\question Let $$A=\{x | x^2-3x+4=0\} \qquad B=\{x \in\mathbb{Z} | x \textrm{ is divisible by 4} \} \qquad C=\{x \in \mathbb{Z} | x \textrm{ is even}\}.$$  Write out the roster notation of each set, using $\dots$ as needed. Which sets are subsets of each other? 

\question Let $A=\{1\}, B=\{\{1\},1\}$ and $C=1.$  Using $A$, $B$ and $C$ and the symbols $\subset, \subseteq$ and $\in,$ write as many true statements as possible.  

\question We will see the following two definitions later in the course, but you are probably already familiar with them.
\begin{definition} An integer $n$ is \emph{even} if $n=2a$ for some integer $a \in \mathbb{Z}$. \end{definition}
\begin{definition} An integer $n$ is \emph{odd} if $n=2a+1$ for some integer $a \in \mathbb{Z}$. \end{definition}
Read the proofs below. What is being proven? Are there any gaps or holes in the logic?  The first one is well-written while the second is not. What are some things that stand out as bad in the second proof?

\begin{proof} Suppose the integer $x$ is odd. Then $x = 2a+1$ for some $a \in \mathbb{Z}$, by definition
of an odd number. Thus $x^2 = (2a+1)^2 = 4a^2+4a+1 = 2(2a^2+2a)+1$, so
$x^2 = 2b +1$ where $b = 2a^2 +2a \in  \mathbb{Z}$. Therefore $x^2$ is odd, by definition
of an odd number. \end{proof}

\begin{proof} $x = 2k + 1.$
Square it so
\[
x^2 = (2k+1)^2 = 4k^2 + 4k + 1.
\]
That's even, plus one. That's odd. {So it works.}
\end{proof}

\question Adapt the good proof above (in problem 9) to prove a complimentary statement about even numbers.

\end{questions}

\end{document}
